% CONCLUSÃO — sem citações novas, conforme NBR 14724

Retomado o objetivo proposto na Introdução, sintetiza-se o rendimento do teste empírico dos cinco vetores de captura. A investigação reuniu, no Capítulo 1, cinco réguas conceituais-normativas que operam em ângulos distintos sobre a transição energética.

No Capítulo 2, essa régua composta foi aplicada à expansão da energia eólica no Ceará e no Rio Grande do Norte entre 2004 e 2025. Demonstra-se a linearidade institucional que liga a doutrina formalizada na década de 1990 ao arranjo eólico contemporâneo.

Sustenta-se, em primeiro achado, que a linearidade institucional opera factualmente, e não por analogia retórica. Em segundo achado, os cinco vetores operam simultaneamente.

Outrossim, em terceiro achado, o discurso de transição verde, quando desprovido de critério de racionalidade ambiental substantiva, configura mecanismo de redistribuição fiscal sem participação política e sem proteção multiespécie.

Propõe-se, como fechamento, que uma transição energética compatível com a racionalidade ambiental exige consulta livre, prévia e informada obrigatória, redistribuição fiscal ancorada em indicadores sociais municipais e monitoramento público independente.
